\subsection{ナイキスト線図}

単位負帰還系の閉ループ伝達関数は
\begin{equation}
  T(s)=\frac{L(s)}{1+L(s)}
\end{equation}
である。したがって閉ループの安定性は、特性方程式
\begin{equation}
  1+L(s)=0
\end{equation}
の根が右半平面に存在するかどうかで決まる。ボード線図の安定余裕は、点 $-1$ への近さを周波数ごとに評価する指標であった。ナイキスト線図はさらに一歩進めて、複素平面上の曲線が点 $-1$ を何回囲むかから、閉ループ右半平面極の個数そのものを数える。

以下では $L(s)$ を実係数の有理関数とし、虚軸上に開ループ極を持たず、調べる閉ループがちょうど安定限界にない、すなわち $1+L(j\omega)\ne0$ であると仮定する。虚軸上に積分器や共振極がある場合には、虚軸を小半円で避ける修正版のナイキスト経路を使う。本書では、右半平面を囲むナイキスト経路を、虚軸上で $-jR$ から $jR$ へ進み、右半平面の大半円で $jR$ から $-jR$ へ戻る向きに取る。この向きは右半平面を時計回りに囲む向きである。

\begin{definition}[囲み数の符号規約]
開ループ $L(s)$ の右半平面極の個数を $P$、閉ループ特性方程式 $1+L(s)=0$ の右半平面根の個数を $Z$ とする。ナイキスト線図が点 $-1$ を時計回りに囲む回数を $N$ と数える。反時計回りの囲みは負として数える。したがって、点 $-1$ を反時計回りに 1 回囲む場合は $N=-1$ である。
\end{definition}

この符号規約の下で、ナイキスト安定判別は
\begin{equation}
  Z=N+P
\end{equation}
と書ける。閉ループ安定性には $Z=0$ が必要十分であるから、開ループが安定で $P=0$ なら $N=0$、すなわち点 $-1$ を囲まないことが必要である。一方、開ループに右半平面極があるときは違う。たとえば $P=1$ なら、閉ループを安定にするには $N=-1$、すなわち点 $-1$ を反時計回りに 1 回囲む必要がある。

ここで数えている $Z$ は抽象的な複素解析上の個数ではなく、時間応答に現れる発散モードの個数である。閉ループ極 $s=\sigma+j\omega$ は過渡応答に $e^{\sigma t}$ と角周波数 $\omega$ の振動成分を与えるため、$Z>0$ ならどこかの初期誤差や外乱成分が時間とともに増幅される。したがってナイキスト線図の囲みは、周波数平面の作図を、時系列での減衰・持続振動・発散へ結び付ける安定性の証拠になる。

\begin{derivation}
偏角原理の考え方を、閉ループ特性関数
\begin{equation}
  F(s)=1+L(s)
\end{equation}
へ適用する。$F(s)$ の零点は閉ループ特性方程式の根であり、$F(s)$ の極は、約分で隠れていない開ループ $L(s)$ の極である。右半平面を時計回りに囲むナイキスト経路を $\Gamma$ とすると、$s$ が $\Gamma$ を一周するときの $F(s)$ の偏角変化は
\begin{equation}
  \Delta_\Gamma \arg F(s)=-2\pi(Z-P)
\end{equation}
である。右半平面内の零点は時計回り経路に対して $-2\pi$ の寄与を持ち、極は符号が反対になるためである。

ここで $F(s)=1+L(s)$ は、$L(s)$ のナイキスト線図を複素平面で右へ 1 だけ平行移動したものである。したがって、$F(s)$ が原点を囲む回数は、$L(s)$ が点 $-1$ を囲む回数と等しい。時計回りの囲み数を
\begin{equation}
  N=-\frac{1}{2\pi}\Delta_\Gamma \arg F(s)
\end{equation}
と定義したので
\begin{equation}
  N=Z-P
\end{equation}
であり、整理すると $Z=N+P$ を得る。
\end{derivation}

実係数の伝達関数では $L(-j\omega)=\overline{L(j\omega)}$ であるため、負の周波数側のナイキスト線図は正の周波数側の鏡映になる。ただし安定判別で数えるのは、正の周波数だけの曲線ではなく、負の周波数、正の周波数、大半円の像を合わせた閉じた曲線である。

\begin{example}[不安定開ループ極を持つナイキスト判別]
開ループを
\begin{equation}
  L(s)=\frac{2}{s-1}
\end{equation}
とする。これは $s=1$ に右半平面極を 1 個持つので
\begin{equation}
  P=1
\end{equation}
である。虚軸上では
\begin{align}
  L(j\omega)
  &=\frac{2}{-1+j\omega}\\
  &=\frac{2(-1-j\omega)}{1+\omega^2}\\
  &=-\frac{2}{1+\omega^2}
    -j\frac{2\omega}{1+\omega^2}
\end{align}
である。$L(j\omega)=x(\omega)+jy(\omega)$ とおくと
\begin{equation}
  x(\omega)=-\frac{2}{1+\omega^2},
  \qquad
  y(\omega)=-\frac{2\omega}{1+\omega^2}
\end{equation}
であり、
\begin{align}
  \{x(\omega)+1\}^2+y(\omega)^2
  &=
  \left(\frac{\omega^2-1}{1+\omega^2}\right)^2
  +
  \left(\frac{2\omega}{1+\omega^2}\right)^2\\
  &=1
\end{align}
を満たす。したがって虚軸部分の像は、中心 $-1$、半径 1 の円である。

$\omega$ を $-\infty$ から $+\infty$ へ増やすと、$L(j\omega)$ は原点近くの上側から出発し、$\omega=0$ で $-2$ を通り、原点近くの下側へ戻る。中心 $-1$ から見た偏角は $0$ から $\pi$ を通って $2\pi$ へ増えるので、点 $-1$ を反時計回りに 1 回囲む。大半円上では $L(s)\to0$ であり、その像は原点近くに縮むため、点 $-1$ の囲み数を変えない。よって本書の符号規約では
\begin{equation}
  N=-1
\end{equation}
である。

ナイキスト判別から
\begin{equation}
  Z=N+P=-1+1=0
\end{equation}
となる。閉ループ右半平面極は存在せず、閉ループは安定である。実際に直接計算しても
\begin{align}
  1+\frac{2}{s-1}&=0,\\
  \frac{s+1}{s-1}&=0
\end{align}
より閉ループ極は $s=-1$ であり、右半平面にはない。

この直接計算は、囲み数の結果を時間応答で読み替えている。閉ループ極が $-1$ にあるため、初期値や外乱で生じた過渡成分は $e^{-t}$ に比例して減衰する。ナイキスト線図が安定と判定したことは、単に曲線が条件を満たすという意味ではなく、この一次モードが右半平面から左半平面へ移り、発散しない応答になったことを表している。

同じ形を
\begin{equation}
  L_K(s)=\frac{K}{s-1}
\end{equation}
として見ると、閉ループ極は $s=1-K$ である。$K=1$ では $s=0$ となり、ナイキスト線図は $\omega=0$ でちょうど点 $-1$ を通る。$K>1$ では点 $-1$ を反時計回りに 1 回囲み、$Z=0$ で安定である。$0<K<1$ では点 $-1$ を囲まず、$N=0$、$Z=1$ となって不安定である。不安定な開ループ極がある場合、「点 $-1$ を囲まない」ことは安定条件ではなく、必要な反時計回りの囲みを失っていることを意味する。
\end{example}

\wfig{frequency-design-examples} 上左は、この不安定開ループ極を持つ一次例でゲインだけを変えたナイキスト線図である。$K=0.6$ の曲線は点 $-1$ を囲まず、$K=1$ では曲線が点 $-1$ を通って安定限界に乗る。$K=2$ では負の周波数から正の周波数へ進む向きで点 $-1$ を反時計回りに 1 回囲むので、本書の符号規約では $N=-1$ となり、$P=1$ と合わせて $Z=0$ が得られる。囲み数は曲線の形だけでなく、開ループ不安定極の個数と向きの規約を同時に読んで決まる。この図の三つの曲線は、同じ数式の時間応答が $K=0.6$ で $e^{0.4t}$ に比例して発散し、$K=1$ で減衰しない境界に乗り、$K=2$ で $e^{-t}$ に比例して減衰することを、周波数応答だけから分けている。

この例は、ゲイン余裕と位相余裕の解釈にも注意を与える。$L(s)=2/(s-1)$ では
\begin{equation}
  \abs{L(j\omega)}=\frac{2}{\sqrt{1+\omega^2}}
\end{equation}
なので、ゲイン交差周波数は
\begin{equation}
  \omega_{gc}=\sqrt{3}
\end{equation}
である。このとき位相は $-120^\circ$ であり、形式的な位相余裕は $60^\circ$ である。一方、位相が $-180^\circ$ となるのは $\omega=0$ で、そのとき $\abs{L(j0)}=2$ である。境界に達するゲイン倍率は $1/2$ であり、この系ではゲインを 2 倍まで増やせるという意味ではなく、ゲインが半分に低下すると安定限界へ落ちるという意味になる。

したがって、この例で位相余裕だけを見ると安定性を取り違える。名目ゲイン $K=2$ からゲインが低下すると閉ループ極 $1-K$ は $-1$ から原点へ近づき、応答は遅くなり、$K=1$ で外乱や初期誤差が減衰しない。さらにゲインが下がると極は右半平面へ入り、時間応答は発散する。安定余裕は「どちら向きのパラメータ変化が危険か」までナイキスト線図と極位置で確認する必要がある。

一般のロバスト設計では、まず $P,N,Z$ により閉ループ安定性を確認し、そのうえでナイキスト線図が点 $-1$ へどれだけ接近するかを見る。感度関数は
\begin{equation}
  S(s)=\frac{1}{1+L(s)}
\end{equation}
であるから、$\abs{1+L(j\omega)}$ が小さい周波数では $\abs{S(j\omega)}$ が大きくなり、モデル誤差や外乱に対する余裕が小さくなる。ゲイン余裕と位相余裕はこの接近の一部を低次元の数値で要約する指標であり、特に不安定開ループ極がある場合には、ナイキスト判別の囲み数を置き換えるものではない。

距離 $\abs{1+L(j\omega)}$ が小さい周波数は、時間応答では減衰の弱い振動や外乱増幅として現れやすい。点 $-1$ の近くを通っても囲み数が変わらなければ名目上の閉ループ安定性は保たれるが、小さなモデル誤差や遅れで曲線が点 $-1$ を横切る余地が残る。この意味で、ナイキスト線図は安定か不安定かだけでなく、安定な応答がどれだけ壊れやすいかも示している。

ナイキスト判別は、むだ時間や高次系を含む場合にも周波数応答から閉ループ右半平面極の個数を評価する。むだ時間 $e^{-Ls}$ はゲインを変えず、位相だけを
\begin{equation}
  \angle e^{-j\omega L}=-\omega L
\end{equation}
だけ遅らせる。高周波ほど位相遅れが増えるため、むだ時間はフィードバックの帯域を強く制限する。

時間領域で見ると、むだ時間は古い誤差に基づいて入力を出すことを意味する。交差周波数付近でこの遅れが大きいと、補正入力が現在の誤差を打ち消す代わりに振動を加える向きへ回り、位相余裕の低下、振動の増大、最終的な不安定化として現れる。

\subsection{根軌跡}

制御器を単純なゲイン $K$ として $L(s)=KG(s)$ と書くと、閉ループ極は
\begin{equation}
  1+KG(s)=0
\end{equation}
を満たす。$K$ を非負実数として変化させたときの閉ループ極の軌跡を根軌跡と呼ぶ。

根軌跡は開ループ極から出発し、開ループ零点または無限遠へ向かう。実軸上では、右側にある実極・実零点の数が奇数の区間に軌跡が存在する。根軌跡上の閉ループ極位置から、指数減衰率および振動周波数を評価する。

具体的には、閉ループ極が $s=\sigma+j\omega$ にあるとき、対応する過渡成分は $e^{\sigma t}\cos\omega t$ と $e^{\sigma t}\sin\omega t$ の形で現れる。根軌跡が左へ動けば減衰は速くなり、虚軸に近づけば減衰は弱くなり、虚軸を越えると振動または偏差が増幅される。したがって根軌跡は、ゲイン変更が時間応答の速さ、振動、安定余裕をどの向きへ変えるかを評価する図である。

二次標準形と対応させると、望ましい極は
\begin{equation}
  s=-\zeta\omega_n\pm j\omega_n\sqrt{1-\zeta^2}
\end{equation}
である。根軌跡がこの近傍を通る場合、ゲイン調整により対応する二次近似の極を実現できる。通らない場合には、零点または極を持つ補償器により根軌跡を変更する。

\wfig{frequency-design-examples} 上右は、三次の標準例
\begin{equation}
  G_3(s)=\frac{1}{s(s+2)(s+4)}
\end{equation}
に対する根軌跡である。閉ループ特性多項式は
\begin{equation}
  s^3+6s^2+8s+K
\end{equation}
となり、Routh--Hurwitz 判別法から安定範囲は $0<K<48$ である。図では $K$ を増やすと二本の枝が実軸上で合流して複素平面へ出て、$K=48$ で虚軸を横切る。したがって、根軌跡は「ゲインを上げると極がどちらへ動くか」を示すだけでなく、ゲイン調整だけで満たせる安定範囲と、補償器で軌跡そのものを変える必要がある範囲を分ける図である。

$K=48$ では補助多項式から虚軸上の極対が現れ、振動が減衰しない境界になる。$K$ をこの値に近づけるほど支配極の実部は 0 に近づき、時間応答では整定が遅く、振動が残りやすい。$K>48$ では根軌跡が右半平面へ入るため、単にオーバシュートが増えるだけでなく、閉ループそのものが不安定になる。

\subsubsection{特性方程式から導く作図規則}

開ループを、互いに共通因子を持たない多項式 $A(s)$、$B(s)$ により
\begin{equation}
  G(s)=\frac{B(s)}{A(s)}
  =\frac{\prod_{j=1}^{m}(s-z_j)}
  {\prod_{i=1}^{n}(s-p_i)}
\end{equation}
と書く。$p_i$ は開ループ極、$z_j$ は開ループ零点であり、$n\ge m$ とする。単位負帰還で比例ゲイン $K$ を掛けると、閉ループ特性方程式は
\begin{equation}
  1+KG(s)=0
\end{equation}
である。分母を払うと
\begin{equation}
  A(s)+K B(s)=0
\end{equation}
となる。したがって根軌跡は、$K\ge0$ に対する多項式 $A(s)+K B(s)$ の根の集合である。$K$ の単位は $KG(s)$ が無次元になるように決まり、以下の幾何的規則では $K$ を正の実数として扱う。

\begin{formula}[根軌跡の角度条件と大きさ条件]
$K>0$ の根軌跡上の点 $s$ は
\begin{align}
  \angle G(s)&=(2\ell+1)\pi,\qquad \ell\in\mathbb{Z},\\
  K&=\frac{1}{\abs{G(s)}}
\end{align}
を満たす。
\end{formula}

\begin{derivation}
$1+KG(s)=0$ は
\begin{equation}
  KG(s)=-1
\end{equation}
と同値である。$K$ は正の実数なので偏角を変えない。右辺 $-1$ の偏角は奇数倍の $\pi$ であるから、根軌跡上では
\begin{equation}
  \angle G(s)=(2\ell+1)\pi
\end{equation}
でなければならない。大きさを比較すると
\begin{equation}
  K\abs{G(s)}=1
\end{equation}
であり、$K=1/\abs{G(s)}$ を得る。
\end{derivation}

角度条件は「その点が閉ループ極になり得るか」を判定し、大きさ条件は「その点を極にするゲインはいくらか」を与える。したがって、根軌跡上の一点を選ぶことは、同時に減衰率 $\sigma$、振動周波数 $\omega$、必要ゲイン $K$ を評価することに等しい。作図規則はこの二つの条件を手計算しやすい形に分解したものである。

この二つの条件から、手作図で使う規則が得られる。まず $K=0$ では特性方程式が $A(s)=0$ となるため、根軌跡は開ループ極 $p_i$ から出発する。$K$ を大きくすると、$m$ 本の枝は有限零点 $z_j$ へ向かい、残りの $n-m$ 本は無限遠へ向かう。したがって根軌跡の枝の本数は開ループ極の個数 $n$ である。

実軸上の点 $\sigma$ を考える。実係数系では、実軸上の点から共役複素極・零点へ引いた角度は互いに打ち消し合う。したがって角度条件に寄与するのは、実軸上にある極・零点だけである。$\sigma$ より右側に存在する実極と実零点の総数を数え、その数が奇数なら
\begin{equation}
  \angle G(\sigma)=(2\ell+1)\pi
\end{equation}
となる。偶数なら角度は偶数倍の $\pi$ となり、$K>0$ の根軌跡には含まれない。

有限零点より開ループ極の方が多い場合、無限遠へ向かう枝は直線漸近線に近づく。$n-m$ 本の漸近線の交点を重心と呼び、
\begin{equation}
  \sigma_a=
  \frac{\sum_{i=1}^{n}p_i-\sum_{j=1}^{m}z_j}{n-m}
\end{equation}
で与えられる。漸近線の角度は
\begin{equation}
  \theta_q=\frac{(2q+1)\pi}{n-m},
  \qquad q=0,1,\ldots,n-m-1
\end{equation}
である。実係数系では複素極と複素零点が共役対で現れるため、重心 $\sigma_a$ は実数になる。

実軸上で二本の枝が合流して複素平面へ出る点を分岐点、複素平面から実軸へ戻る点を合流点と呼ぶ。実軸上で
\begin{equation}
  K(s)=-\frac{A(s)}{B(s)}
\end{equation}
とおくと、分岐点または合流点の候補は
\begin{equation}
  \frac{\dd K}{\dd s}=0
\end{equation}
を満たす。すなわち
\begin{equation}
  A'(s)B(s)-A(s)B'(s)=0
\end{equation}
である。ただし、この方程式の解がすべて根軌跡上の分岐点になるわけではない。候補点が実軸上の根軌跡区間にあり、さらに $K(s)>0$ を満たすことを確認する必要がある。

根軌跡が虚軸を横切るゲインは、閉ループ特性多項式 $A(s)+KB(s)$ に Routh--Hurwitz 判別法を適用して求める。$K$ を含む Routh 表を作り、第 1 列の符号が変わる値または行が 0 になる値を調べる。行が 0 になる場合には、その一つ上の行から補助多項式を作り、$s=j\omega$ 上の交点を求める。この計算により、根軌跡の図だけでは読みにくい安定限界ゲインを代数的に得られる。

複素極から根軌跡が出発する角度も角度条件から決まる。複素極 $p_k$ の近くで
\begin{equation}
  s=p_k+\rho e^{j\theta_d},
  \qquad \rho>0
\end{equation}
とおくと、$s-p_k$ の偏角だけが未知の出発角 $\theta_d$ である。$\rho\to0$ の極限で角度条件を使うと
\begin{equation}
  \theta_d
  =\pi
  +\sum_{j=1}^{m}\angle(p_k-z_j)
  -\sum_{\substack{i=1\\i\ne k}}^{n}\angle(p_k-p_i)
  \pmod{2\pi}
\end{equation}
を得る。角度は正の実軸から反時計回りを正として測る。共役な下側の極では、実係数系なら出発角も共役対称になる。

\begin{example}[複素極を含む四次系の根軌跡]
無次元化された開ループ
\begin{equation}
  G(s)=\frac{1}{s(s+4)(s^2+2s+5)}
\end{equation}
を考える。実機のゲインには単位が付くが、ここでは根軌跡の幾何を示すため、時間とゲインを正規化して扱う。開ループ極は
\begin{equation}
  0,\qquad -4,\qquad -1+2j,\qquad -1-2j
\end{equation}
であり、有限零点はない。閉ループ特性方程式は
\begin{align}
  1+KG(s)&=0,\\
  s(s+4)(s^2+2s+5)+K&=0
\end{align}
である。多項式を展開すると
\begin{equation}
  s^4+6s^3+13s^2+20s+K=0
\end{equation}
となる。

実軸上では、区間 $(-4,0)$ にある点の右側には実極 $0$ が 1 個あるため、この区間に根軌跡が存在する。$0$ より右側では右側の実極・実零点の数が 0 個、$-4$ より左側では 2 個なので、どちらも根軌跡ではない。

有限零点がないので、4 本すべての枝が無限遠へ向かう。漸近線の重心は
\begin{equation}
  \sigma_a=
  \frac{0+(-4)+(-1+2j)+(-1-2j)}{4}
  =-\frac{3}{2}
\end{equation}
である。漸近線の角度は
\begin{equation}
  45^\circ,\quad 135^\circ,\quad 225^\circ,\quad 315^\circ
\end{equation}
である。したがって十分大きい $K$ では、二本の枝は右半平面へ向かう漸近線に近づく。

分岐点候補は
\begin{equation}
  K(s)=-(s^4+6s^3+13s^2+20s)
\end{equation}
の停留点から求める。
\begin{align}
  \frac{\dd K}{\dd s}
  &=-(4s^3+18s^2+26s+20)=0,\\
  4s^3+18s^2+26s+20&=0
\end{align}
である。この三次方程式の実根のうち、根軌跡区間 $(-4,0)$ にあるものは
\begin{equation}
  s_b\simeq -2.826
\end{equation}
であり、このとき
\begin{equation}
  K_b=K(s_b)\simeq 24.33
\end{equation}
である。したがって実極 $0$ と $-4$ から出発した二本の枝は、$K\simeq24.33$ で実軸上の $s\simeq-2.826$ に集まってから複素平面へ出る。

虚軸交差を求めるため、特性多項式
\begin{equation}
  s^4+6s^3+13s^2+20s+K
\end{equation}
の Routh 表を作る。
\begin{equation}
\begin{array}{c|ccc}
s^4 & 1 & 13 & K\\
s^3 & 6 & 20 & 0\\
s^2 & 29/3 & K & 0\\
s^1 & 20-\dfrac{18K}{29} & 0 & 0\\
s^0 & K & 0 & 0
\end{array}
\end{equation}
第 1 列が正である条件は
\begin{equation}
  K>0,\qquad 20-\frac{18K}{29}>0
\end{equation}
である。よって閉ループは
\begin{equation}
  0<K<\frac{290}{9}
\end{equation}
で安定であり、$K=290/9$ で虚軸を横切る。このとき $s^1$ 行が 0 になるので、$s^2$ 行から補助多項式
\begin{equation}
  \frac{29}{3}s^2+K=0
\end{equation}
を作る。$K=290/9$ を代入すると
\begin{equation}
  s^2=-\frac{10}{3}
\end{equation}
であり、交点は
\begin{equation}
  s=\pm j\sqrt{\frac{10}{3}}
\end{equation}
である。

この交点では、閉ループ応答に角周波数 $\sqrt{10/3}$ の減衰しない振動成分が現れる。$K$ が $290/9$ より少し小さければ対応する極対は左半平面にあり、振動は時間とともに減衰する。$K$ がこの値を越えると極対は右半平面へ移り、同じ周波数付近の振動包絡が増加する。Routh 表の第 1 列は、この時間応答の変化が起こるゲインを図からの目測ではなく代数的に固定している。

最後に、上側の複素極 $p_k=-1+2j$ からの出発角を求める。零点はない。ほかの極から $p_k$ へ向かうベクトルの角度は
\begin{align}
  \angle\{(-1+2j)-0\}&\simeq116.6^\circ,\\
  \angle\{(-1+2j)-(-4)\}&\simeq33.7^\circ,\\
  \angle\{(-1+2j)-(-1-2j)\}&=90^\circ
\end{align}
である。これらの和は約 $240.3^\circ$ であるから
\begin{equation}
  \theta_d
  =180^\circ-240.3^\circ
  \simeq -60.3^\circ
\end{equation}
となる。上側の複素極から出る枝は正の実軸から見て約 $-60^\circ$ の方向へ出発し、下側の複素極から出る枝はその共役対称として約 $+60^\circ$ の方向へ出発する。

この例では、$K=0$ の極配置、実軸区間、漸近線、分岐点、虚軸交差、複素極の出発角がすべて同じ特性方程式から計算できる。出発角は、複素極から出る枝が初期に右向き成分を持ち、ゲイン増加で減衰が弱くなり得ることを示す。分岐点 $s_b$ は実極由来の二つのモードが同じ減衰率で合流してから振動モードへ変わる位置であり、虚軸交差はその振動モードが減衰を失う限界である。特に Routh 表から $K=290/9$ を越えると閉ループ極が右半平面へ入ることが分かるため、根軌跡の作図結果は安定なゲイン範囲の計算と時間応答の解釈を合わせて評価する。
\end{example}
